\begin{lemma}
  Let $M$ be a smooth manifold with affine connection $\nabla$. If $h \colon \mathbb{R} \to M$ is a smooth curve, then
  \begin{equation}
    e
  \end{equation}
\end{lemma}
\begin{proof}
Let $p \in \mathbb{R}$. There exists an open set $U \subseteq M$  containing $h(p)$ on which there exists a moving frame $E_1, \ldots, E_n \colon U \to \mathrm{T}M|_U$. In this frame, the connection is expressible in terms of a linear operator. Explicitly, let $I = h^{-1}(U)$ and $\gamma = h|_I$, and suppose that $V \colon I \to \gamma^*(\mathrm{T}M|_U) \cong I \times \mathbb{R}^n$ is a vector field along $\gamma$. Then
\begin{equation}
(\gamma^* \nabla)_{\mathrm{d}/\mathrm{d}t} V = \der{V}{t} + \omega V
\end{equation}
where $\omega(t) \colon \mathrm{T}_{\gamma(t)}M \to \mathrm{T}_{\gamma(t)}M $ is the endomorphism defined by $\omega E_i = \nabla_{\mathrm{d}/\mathrm{d}t} E_i$. 

\end{proof}

\newpage

\begin{lemma}
  Let $\nabla$ be an affine connection on $\mathbb{R}^n$ and denote by $\omega$ its connection 1-form with respect to the standard frame. For any smooth family of curves $\gamma \colon A \times I \to \mathbb{R}^n$, there exists $\mathcal{E} \colon A \times \mathbb{R}^n \to \mathbb{R}^n$ such that
  \begin{equation}
    \Gamma_\nabla (\gamma_\alpha) = \Big[\mathbf{1} - \, \omega (\gamma'_\alpha(0))\Big] + \mathcal{E}_\alpha(\gamma_\alpha'(0))
  \end{equation}
  with $\mathcal{E}_\alpha(u) = O_\alpha(|u|^3)$ as $u \to 0$. 
\end{lemma}


\begin{lemma}
  Let $\nabla$ be an affine connection on $\mathbb{R}^n$ and denote by $\omega$ its connection 1-form with respect to the standard frame. For any smooth family of curves $\gamma \colon A \times I \to \mathbb{R}^n$, there exists $\mathcal{E} \colon A \times \mathbb{R}^n \to \mathbb{R}^n$ such that
  \begin{equation}
    \tau_\nabla (\gamma_\alpha) = \Big[\mathbf{1} + \tfrac{1}{2} \, \omega (\gamma'_\alpha(0))\Big]  \gamma'_\alpha(0) + \mathcal{E}_\alpha(\gamma_\alpha'(0))
  \end{equation}
  with $\mathcal{E}_\alpha(u) = O_\alpha(|u|^3)$ as $u \to 0$. 
\end{lemma}

\begin{proof}

\end{proof}

\begin{lemma}
  Let $U \subseteq \mathbb{R}^n$ be open, $\nabla$ be an affine connection on $U$, and $A$ be a smooth manifold. For a smooth family of smooth curves $\{\gamma_\alpha \colon I \to U\}_{\alpha \in A}$, we have
  \begin{equation}
    \tau_\nabla\Big[\gamma_\alpha(0 \to t)\Big] = t \cdot \gamma_\alpha'(0) + O_\alpha(t^2).
  \end{equation}
\end{lemma}
\begin{proof}
For each $\alpha \in A$, let $\omega^\alpha =  \omega^\alpha_1\, \d t \in \Omega^1(I, \mathfrak{gl}_n)$ be the connection $1$-form associated to $\gamma_\alpha^* \nabla$. Let $J \subseteq I$ be some bounded interval containing $0$. By Taylor's theorem, for all $t \in J$, we have 
\begin{equation}
  \Gamma(\gamma_\alpha)_{t}^{0} \left[\gamma_\alpha'(t)\right] = \Big[\mathbf{1} + \mathcal{E}_\alpha(t)\Big] \gamma_\alpha'(t),
\end{equation}
where $\lVert\mathcal{E}_\alpha(t)\rVert_{\max} \leq \sup_{J} \left\lVert \omega^\alpha_1 \right\rVert _{\max} \cdot t$. Integrating, we have for all $t \in J$,
\begin{equation}
  \tau_\nabla\Big[\gamma_\alpha(0 \to t)\Big] = \big[\gamma_\alpha(t) - \gamma_\alpha(0)\big] + \epsilon_\alpha(t),
\end{equation}
where 
\begin{equation}
|\epsilon_\alpha(t)|_{\max} \leq \int_{0}^{1} |\mathcal{E}_\alpha(t) \gamma_\alpha'(t)|_{\max} \, \d t \leq  n \cdot \sup_{J}{\lVert \omega^\alpha_1 \rVert_{\max}} \cdot \sup_{J} |\gamma_\alpha'| \cdot \tfrac{1}{2}t^2.
\end{equation}
Since $\omega_1$ and $\gamma'$ are smooth in $\alpha$, both supremums are continuous in $\alpha$, so $\epsilon_\alpha(t) = O_\alpha(t^2)$. Applying Taylor's theorem again, for all $t \in J$ we may write
\begin{equation}
  \gamma_\alpha(t) - \gamma_\alpha(0) = t \cdot \gamma_\alpha'(0) + \delta_\alpha(t),
\end{equation}
where $|\delta_\alpha(t)|_{\max} \leq \sup_J |\gamma_\alpha''| \cdot \tfrac{1}{2} t^2$. Since $\gamma''$ is smooth in $\alpha$, it follows that $\delta_\alpha(t) = O_\alpha(t^2)$ as well. The conclusion follows.
\end{proof}


\newpage

\begin{theorem}
Let $\nabla$ be an affine connection on $\mathbb{R}^n$. For all vectors $u, v \in \mathbb{R}^n$, we have
\begin{equation}
  \lim_{x,y \to 0}\, \mfrac{1}{xy} \cdot \tau_{\nabla} \bigl\llbracket x \cdot u, y \cdot v \bigr\rrbracket{}_p  = T_p(u, v).
\end{equation}
uniformly on compact subsets.
\end{theorem}

\begin{proof}
  For brevity, denote by $\pi_{p, x, y} = \bigl\llbracket x \cdot u, y \cdot v \bigr\rrbracket{}_p$ the parallelogram path at $p$ with side lengths $x$ and $y$. Let $A \subset \mathbb{R}^n$ be a compact set and let $Q \subseteq \mathbb{R}^2$ be a compact neighborhood of the origin. Now, for integers $0 \leq k < 4$,  define 
  \begin{align}
    \gamma^{k}_{p, x, y} = \pi_{p, x, y}(k \to k + 1), \quad \quad q^{k}_{p, x,y} = \gamma^k_{p, x, y}(0), \quad \quad  e^k_{p, x, y} = {\gamma^k_{p, x, y}}{}'(0) 
  \end{align}
  to be the sides, vertices, and tangent vectors along $\pi$. Note that $e^k$ does not depend on $p$ as a vector in $\mathbb{R}^n$, so we may write $e^k_{p, x, y} = e^k_{x,y}$ when the tangent space the vector lies in is not relevant.
  By Proposition \ref{}, $\tau_\nabla(\pi)$ is smooth as a function of $p$ and $(x,y)$. In particular, we have
  \begin{equation}
    \tau_\nabla(\pi_{p, x, y}) = xy \cdot \partial_{xy} \tau_{\nabla}(\pi)|_{p, 0, 0} + O_p(x^2y + xy^2).
  \end{equation}
  by Taylor's theorem, since $\tau_\nabla(\pi_{p,x,0}) = 0$ and $\tau_\nabla(\pi_{p,0,y}) = 0$

  By Proposition \ref{}, the parallel transport along the reverse of each curve may be written as
  \begin{equation}
    \Gamma \left(\overline{\gamma}{}^k_{p,x,y}\right)  = \Big[\mathbf{1} + \omega(e^k_{p, x, y})\Big] + \mathcal{E}(q^k_{p,x,y},\, e^k_{p,x, y})
  \end{equation}
  where $\mathcal{E}(q, v) = O_q(v^2)$ as $v \to 0$. Likewise, by Proposition \ref{}, the accumulated torsion along each curve may be written as
  \begin{align}
    \tau_\nabla(\gamma^k_{p, x, y}) = e^k_{x, y} + \epsilon_{p}(|e^k_{x, y}|),
  \end{align} 
  where $\epsilon_p(v) = O_p(v^2)$ as $t \to 0$. Therefore, for $1 \leq k \leq 4$, we may write
  \begin{equation}
    \left[\prod_{i < k} \Gamma^i_{x, y}(t_i)\right] \tau^k_{x,y}(t_k) = e^k_{x, y} +\sum_{i < k} \big[ \omega(e^i_{p, x, y})\big]e^k_{x, y} + \sum_{i,j,k} O_{p}(x^2y + xy^2). \;
  \end{equation}
  Thus,
  \begin{equation}
    \sum_{k \leq 4}\left[\prod_{i < k} \Gamma^i_{x, y}(t_i)\right] \tau^k_{x,y}(t_k) =  \sum_{k \leq 4} t_k e_k + \sum_{k \leq 4} \tfrac{1}{2} t_k^2 \big[\omega_{p^k_{x,y}}(e_k)\big]e_k + \sum_{\substack{k \mskip 1mu \leq \mskip 1mu 4 \\ i \mskip 2mu < \mskip 1mu k}} t_{k} t_i \big[ \omega_{p^i_{x,y}}(e_i)\big]e_k + \mathcal{E}_{x,y}(t), \;
  \end{equation}
  where $\mathcal{E} \colon \mathbb{R}^2 \times I^4 \to \mathbb{R}^n$ is the sum of the individual error terms. It follows that there exists a neighborhood $U$ of the origin $\mathbf{0} \in I^4$ and a continuous map $C \colon \mathbb{R}^2 \to \mathbb{R}$ such that for all $t \in U$, we have
  \begin{equation}
    |\mathcal{E}_{x,y}(t)|_{\max} \leq C(x,y)  \sum_{i, j, k} |t_i t_j t_k|.
  \end{equation}
  Take $t = (x, y, x, y)$. Since $e_1 = -e_3$ and $e_2 = -e_4$, the first sum vanishes and the second sum cancels out with part of the third, so
  \begin{align}
    \mfrac{1}{xy} \cdot \tau_{\nabla} \bigl\llbracket x \cdot u, y \cdot v \bigr\rrbracket{}_p &= \Big(\big[ \omega_{p^1_{x,y}}(u)\big]v - \big[ \omega_{p^2_{x,y}}(v)\big]u \Big) + \widetilde{\mathcal{E}}(x,y),
  \end{align}
  where $\widetilde{\mathcal{E}}(x,y) = \frac{1}{xy} \cdot \mathcal{E}_{x,y}(x,y,x,y) \leq C_2(x, y) (|x| + |y|)$

\end{proof}

\newpage 
\begin{theorem}
  Let $\nabla$ be an affine connection on $\mathbb{R}^n$ and denote by $\omega$ its connection $1$-form in the standard frame. For all points $p \in \mathbb{R}^n$ and vectors $u, v \in \mathbb{R}^n$, we have 
  \begin{equation}
    \frac{\partial^2}{\partial x \partial y} \Big( \Gamma \bigl\llbracket x \cdot u, y \cdot v \bigr\rrbracket{}_p \Big)_{(0, 0)} =  \big[\mathrm{d}\omega + \omega \wedge \omega\big]{}_p(u, v).
  \end{equation}
\end{theorem}

\begin{proof}
For brevity, let $\pi_{x, y} = \bigl\llbracket x \cdot u, y \cdot v \bigr\rrbracket{}_p$. Consider the ideal 
\begin{equation}
J = \langle x^2, y^2\rangle \cdot \mathbf{1} \subseteq \mathscr{C}^\infty(\mathbb{R}^2, \mathrm{Mat}_{n \times n}(\mathbb{R})),
\end{equation}
which lies in the kernel of the differential operator $\partial_{xy}|_{0,0}$. For $0 \leq k < 4$,  define 
  \begin{align}
    \gamma^{k}_{x, y} = \pi_{x, y}(k \to k + 1), \quad \quad p^{k}_{x,y} = \gamma^k_{x, y}(0), \quad \quad  e^k_{x, y} = {\gamma^k_{x, y}}{}'(0) 
  \end{align}
  to be the sides, vertices, and tangent vectors along $\pi$. By Taylor's theorem, there exist smooth error terms $\mathcal{E}^k \colon \mathbb{R}^2 \to \mathrm{Mat}_{n \times n}(\mathbb{R})$ such that
  \begin{align}
    \Gamma(\gamma^{1}_{x,y}) &= \mathbf{1} - x \cdot \partial_x \Gamma(\gamma^{1}_{x,y})|_{(0, y)} + x^2 \cdot \mathcal{E}^1(x, y), \\
    \Gamma(\gamma^{2}_{x,y}) &= \mathbf{1} - y \hspace{1pt} \cdot \partial_y \Gamma(\gamma^{2}_{x,y})|_{(x, 0)} + y^2 \cdot \mathcal{E}^2(x, y), \\
    \Gamma(\gamma^{3}_{x,y}) &= \mathbf{1} + x \cdot \partial_x \Gamma(\gamma^{3}_{x,y})|_{(0, y)} + x^2 \cdot \mathcal{E}^3(x, y), \\
    \Gamma(\gamma^{4}_{x,y}) &= \mathbf{1} + y \hspace{1pt} \cdot \partial_y \Gamma(\gamma^{4}_{x,y})|_{(x, 0)} + y^2 \cdot \mathcal{E}^4(x, y).
  \end{align}
  By Proposition \ref{}, the parallel transport around $\pi_{x,y}$ is
  \begin{align}
    \Gamma(\pi_{x, y}) &\equiv_J \Big(\mathbf{1} - x \cdot  \omega_{p^1_{x, y}}(u)\Big) \Big(\mathbf{1} - y \cdot  \omega_{p^2_{x, y}}(v)\Big) \Big(\mathbf{1} + x \cdot  \omega_{p^3_{x, y}}(u)\Big) \Big(\mathbf{1} + y \cdot  \omega_{p^4_{x, y}}(v)\Big)
  \end{align}
  Expanding modulo $J$, we obtain
  \begin{align}
    \Gamma(\pi_{x,y}) \equiv_J \mathbf{1} \mskip 3mu &+ \mskip 5mu x \mskip 3mu \cdot \Big(\omega_{p^3_{x,y}}(u) - \omega_{p^1_{x,y}}(u)\Big) + y \cdot \Big(\omega_{p^4_{x,y}}(v) - \omega_{p^2_{x,y}}(v)\Big) \nonumber \\
    &+ xy \cdot \Big( \omega_{p^1_{x,y}}(u)  \omega_{p^2_{x,y}}(v) - \omega_{p^2_{x,y}}(v)  \omega_{p^3_{x,y}}(u) \Big) \nonumber \\
    &- xy \cdot \Big(\omega_{p^1_{x,y}}(u)  - \omega_{p^3_{x,y}}(u) \Big) \omega_{p^{4}_{x,y}}(v).
  \end{align}
  It follows that
  \begin{align}
    \partial_{xy} \Gamma(\pi_{x,y})|_{0, 0} &= \lim_{y \to 0} \mfrac{1}{y} \cdot  \Big[ \omega_{p^3_{0,y}}(u) -  \omega_{p^1_{0,y}}(u) \Big] + \lim_{x \to 0} \mfrac{1}{x} \cdot \Big[ \omega_{p^4_{x,0}}(v) -  \omega_{p^2_{x,0}}(v) \Big] \nonumber \\
    &+ \Big[ \omega_{p^{1}_{0,0}}(u)\omega_{p^{2}_{0,0}}(v) -  \omega_{p^{2}_{0,0}}(u)\omega_{p^{3}_{0,0}}(v) \Big],
  \end{align}
  Evaluating the limits yields $\partial_{xy} \Gamma(\pi_{x,y})|_{0, 0} = \big[\mathrm{d}\omega + \omega \wedge \omega\big]{}_p(u, v)$, as desired. 
\end{proof}

\newpage 
\begin{theorem}
  Let $\nabla$ be an affine connection on $\mathbb{R}^n$. If $\gamma \colon I \to \mathbb{R}^n$ is a smooth curve and $Z \in \Gamma(\mathrm{T}\mathbb{R}^n)$ is a smooth vector field, then
  \begin{equation}
    \der{}{t} \Big[\Gamma_{0}^{t} (\gamma) Z_{\gamma(0)}\Big]_{t = 0} = \left.\der{Z}{t}\right|_{t = 0} - \nabla_{\gamma'(0)} Z.
  \end{equation}
\end{theorem}
\begin{proof}
By Proposition \ref{}, and the Leibniz rule, we have
\begin{equation}
  \nabla_{\gamma'(0)} Z = \der{}{t} \Big[\Gamma_t^0(\gamma) Z_{\gamma(t)} \Big]_{t = 0} = \der{}{t} \Big[\Gamma_t^0(\gamma)\Big]_{t = 0} Z_{\gamma(0)}  + \left.\der{Z}{t}\right|_{t = 0}.
\end{equation}
Since $\Gamma_0^0(\gamma) = \mathbf{1}$, we also have $\der{}{t}\big[\Gamma_t^0(\gamma)\big]_{t = 0} = -\der{}{t}\big[\Gamma_0^t(\gamma)\big]_{t = 0}$, implying the conclusion.
\end{proof}

This may also be rephrased in terms of linear operators on $\Gamma(T\mathbb{R}^n)$. Let $\nabla$ be an affine connection on a smooth manifold $M$, and $X \in \Gamma(\mathrm{T}M)$ be a vector field with flow $\Phi^X \colon I \times M \to M$. Then for each $s, t \in I$ we may define a vector bundle isomorphism
\begin{equation}
  \mathrm{P}(\Phi_X)_s^t \colon \mathrm{T}M \to \mathrm{T}M; \quad \quad \big[\mathrm{P}(\Phi_X)^{\mskip 0.5mu t}_{s}\big] \xi_{p} = \Big[\mathrm{P}(\Phi_X(-, p))_{0}^{t-s}\Big] \xi_{p}. 
\end{equation}
covering the flow diffeomorphism $\Phi_{X}^{\mskip 2mu t-s}$.
By taking the induced map on global sections, we also obtain  $\mathbb{R}$-linear operators
\begin{equation}
\mathrm{P}_\ast(\Phi_X)_{s}^{\mskip 1mu t} \colon \Gamma(\mathrm{T}M) \to \Gamma(\mathrm{T}M); \quad \quad \big[\mathrm{P}_\ast(\Phi_X)_{s}^{\mskip 0.5mu t} Z\big]{}_{p} = \Big[ \mathrm{P}\big(\Phi_X (-, p)\big){}_{0}^{t - s} \Big] Z_{p}.
\end{equation}
The advantage here is that we may now parallel transport a vector field along an entire family of curves, and compose these operators without having to keep track of individual endpoints. Intuitively, the first map is \say{dynamic}, in the sense that it observes the result of translating a vector $\xi_p$ along $\Phi_X$, moving its basepoint along the flow lines. On the other hand, the second map operates on vector fields, which can always be evaluated at the same point in $M$. As such, it can observe the parallel transport along $\Phi_X$ in a \say{static} manner.


\begin{theorem}
  Let $\nabla$ be an affine connection on $\mathbb{R}^n$ and denote  by $\omega$ its connection form in the standard frame. If $X \in \Gamma(\mathrm{T}\mathbb{R}^n)$ is a vector field with global flow $\Phi_X \colon I \times \mathbb{R}^n \to \mathbb{R}^n$, then
  \begin{equation}
    \der{}{t} \Big[ \mathrm{P}(\Phi_X)_{0}^{t}
    \Big]_{t=0} = -\omega(X). 
  \end{equation}
\end{theorem}
\begin{proof}
  Let $v_p \in \mathrm{T}\mathbb{R}$ be a tangent vector based at $p \in \mathbb{R}$. Then $V(t) = \big[\mathrm{P}(\Phi_X)^t_0\big] v_p$ is a parallel vector field along the curve $\gamma(t) = \Phi_X(t, p)$, so it satisfies the differential equation $V' + \omega(\gamma')V = 0$. Since $\gamma'(0) = X_p$ by definition of $\Phi_X$, we have 
  \begin{equation}
    \left.\der{}{t} \Big[\mathrm{P}(\Phi_X)_{0}^{t}\Big]v_p\right|_{t=0} = -\omega(X_p)v_p.
  \end{equation}
  The conclusion follows.
\end{proof}
There is a similar result for the induced map on global sections. Roughly speaking, since $\omega(X) = \nabla_X - D_X$, and $P_\ast$ describes parallel transport in a \say{static} manner,  we expect differentiation to recover only the $\nabla$ term.

\begin{theorem}
  Let $\nabla$ be an affine connection on $\mathbb{R}^n$. If $X \in \Gamma(\mathrm{T}\mathbb{R}^n)$ is a vector field with global flow $\Phi_X \colon I \times \mathbb{R}^n \to \mathbb{R}^n$, then
  \begin{equation}
    \der{}{t} \Big[ \mathrm{P}_\ast(\Phi_X)_{0}^{t}
    \Big]_{t=0} = -\nabla_X. 
  \end{equation}
\end{theorem}
\begin{proof}
Let $V \in \Gamma(\mathrm{T}\mathbb{R}^n)$ be a vector field and $p \in \mathbb{R}^n$. Then 
\begin{equation}
   \Big[\mathrm{P}_*(\Phi_X)_{0}^{t} V\Big]_{\Phi_X(t, p)} = \mathrm{P}(\Phi_X)_{0}^{t} V_p.
\end{equation}
which recovers the original \say{dynamic} description of parallel transport. Now, recall the following fact: for a smooth map $W \colon I \to \Gamma(\mathrm{T}\mathbb{R})$ and curve $\gamma \colon I \to \mathbb{R}^n$, we have
\begin{equation}
  \der{}{t} \Big[W(t)_{\gamma(t)}\Big] = W'(t)_{\gamma(t)} + \big[D_{\gamma'(t)} W(t)\big]{}_{\gamma(t)}.
\end{equation}
Thus, by the previous fact and Proposition \ref{},
\begin{equation}
  \left.\der{}{t}\Big[\mathrm{P}_*(\Phi_X)_{0}^{t} V\Big]_{\Phi_X(t, p)}\right|_{t = 0} = \left[\der{}{t} \mathrm{P}_*(\Phi_X)_{0}^{t} V\Big|_{t = 0} + D_{X} V \right]_{\raisebox{2pt}{$\scriptstyle p$}} = -\omega(X_p)V_p.
\end{equation}
Since $\omega = \nabla_X - D_X$, the conclusion follows.
\end{proof}


\section{Non-abelian Stokes theorems}

The Stokes theorem describes a fundamental...


In each case, we wish to study a certain \say{integral} quantity $A$ which may be defined on smooth oriented $1$-cells (i.e., paths) on a manifold. The Stokes' theorem roughly states that computing this quantity over the boundary of an \say{infinitesimal} oriented $2$-cell yields a differential $2$-form $\alpha$, which when integrated over the interior of a sufficiently nice $2$-cell $c$, reproduces the value of $A$ on the boundary $\partial c$. The main idea behind the theorem is that the $2$-cell $c$ may be decomposed into smaller cells whose boundaries \say{cancel} each other out.

When $A$ lies in an abelian Lie group, this produces the usual Stokes theorem; for completeness, this is the first case we will discuss. However, in the other cases we wish to explore, $A$ is  

\begin{center}
\begin{tabular}{c|c|cc}
$A$ & composition law  \\ \hline
$I(\gamma) = \int_{\gamma} \alpha$ & $\delta f (\gamma_1 \ast \gamma_2) = \delta f(\gamma_1) + \delta f (\gamma_2)$ & \say{abelianizing} groupoid homomorphism  \\[1ex]
$I(\gamma) = P ^\nabla(\gamma) = $ & $P^\nabla(\gamma_1 \ast \gamma_2) = P^\nabla (\gamma_1) P^\nabla (\gamma_2)$ &  groupoid homomorphism \\[1ex] $\tau^\nabla(\gamma)$ & $A(\gamma_1) + B(\gamma_1)A(\gamma_2)$ & crossed groupoid homomorphism
\end{tabular}
\end{center}

\begin{enumerate}
  \item The most basic data that can be integrated over a curve $\gamma \colon I \to M$ is a $1$-form $\alpha$ which produces an integral of a scalar quantity by \say{eating} each tangent vector. Specifically, we define
  \begin{equation}
    \int_{\gamma} \alpha = \int_{I} \alpha(\gamma'(t)) \, \d t.
  \end{equation}
  From this, we obtain the composition law $\int_{\gamma_1 \ast \gamma_2} \alpha = \int_{\gamma_1} \alpha + \int_{\gamma_2} \alpha$.

  \item Given a connection $\nabla$ on $M$, we may identify nearby tangent spaces, providing another means of integrating data over a curve $\gamma \colon I \to M $. Indeed, let $\omega$ be the associated connection form with respect to some frame $E$. At each $t \in I$, the endomorphism $\omega(\gamma'(t))$ of $\mathrm{T}_{\gamma(t)} M$ measures the infinitesimal deviation of the frame from being parallel along the curve. These may be integrated to obtain the parallel transport map along $\gamma$; indeed, by Proposition \ref{}, we have
  \begin{equation}
    \big[P^\nabla(\gamma)\big]_{E} = -\int_I \Big[\omega(\gamma'(t)) \circ P^\nabla(\gamma)_{0}^{t}\Big]_E \, \d t.
  \end{equation}
  While this expresses the parallel transport as an integral, note that $P^\nabla(\gamma)$ appears on both sides, as the individual endomorphisms $\omega(\gamma'(u))$ can only be compared by using parallel transport.

  \item Given a connection $\nabla$ on $M$, we may identify nearby tangent spaces, we may instead integrate the tangent vectors along a curve $\gamma \colon I \to M$, using parallel transport to identify different tangent spaces. This is just the development along $\gamma$, 
  \begin{equation}
    D^\nabla(\gamma) = \int_I \Big[ P^\nabla(\gamma)_{\mskip 1mu t}^{0} \Big]\gamma'(u) \, \d t.
  \end{equation}
  The development satisfies the \say{crossed homomorphism} law $D^\nabla(\gamma_1 \ast \gamma_2) = $.
\end{enumerate}
